
\documentclass{article}
\usepackage{colortbl}
\usepackage{makecell}
\usepackage{multirow}
\usepackage{supertabular}

\begin{document}

\newcounter{utterance}

\centering \large Interaction Transcript for game `hot\_air\_balloon', experiment `air\_balloon\_survival\_en\_complexity\_easy', episode 4 with gpt{-}5{-}mini{-}2025{-}08{-}07{-}t1.0.
\vspace{24pt}

{ \footnotesize  \setcounter{utterance}{1}
\setlength{\tabcolsep}{0pt}
\begin{supertabular}{c@{$\;$}|p{.15\linewidth}@{}p{.15\linewidth}p{.15\linewidth}p{.15\linewidth}p{.15\linewidth}p{.15\linewidth}}
    \# & $\;$A & \multicolumn{4}{c}{Game Master} & $\;\:$B\\
    \hline

    \theutterance \stepcounter{utterance}  
    & & \multicolumn{4}{p{0.6\linewidth}}{
        \cellcolor[rgb]{0.9,0.9,0.9}{
            \makecell[{{p{\linewidth}}}]{
                \texttt{\tiny{[P1$\langle$GM]}}
                \texttt{Stai partecipando a un gioco di negoziazione collaborativa.} \\
\\ 
\texttt{Insieme a un altro partecipante, dovete accordarvi su un unico set di oggetti da conservare. Ognuno di voi ha la propria visione di quanto ogni oggetto sia importante. Non sai come l'altro partecipante valuta gli oggetti. Inoltre, ti viene fornito lo sforzo che ogni oggetto richiede.} \\
\texttt{Potete accordarvi su un set solo se lo sforzo totale degli oggetti selezionati non supera un limite condiviso:} \\
\\ 
\texttt{LIMITE = 9612} \\
\\ 
\texttt{Ecco i valori di sforzo dei singoli oggetti:} \\
\\ 
\texttt{Sforzo oggetto = \{"B21": 935, "B00": 928, "A08": 725, "A27": 218, "B94": 839, "B53": 970, "A23": 394, "C65": 677, "A89": 946, "B67": 530, "B45": 138, "C34": 834, "C12": 734, "A05": 593, "B60": 261, "C92": 743, "B69": 4, "A84": 735, "C06": 124, "B12": 829, "A16": 207, "B85": 782, "C35": 577, "B10": 388, "C15": 679, "B09": 494, "B16": 559, "C81": 629, "A58": 237, "C24": 275, "A30": 39, "A66": 654, "B75": 172, "B46": 688, "C70": 687\}} \\
\\ 
\texttt{Ecco la tua visione personale sull'importanza di ogni oggetto:} \\
\\ 
\texttt{Valori di importanza dell'oggetto = \{"B21": 138, "B00": 583, "A08": 868, "A27": 822, "B94": 783, "B53": 65, "A23": 262, "C65": 121, "A89": 508, "B67": 780, "B45": 461, "C34": 484, "C12": 668, "A05": 389, "B60": 808, "C92": 215, "B69": 97, "A84": 500, "C06": 30, "B12": 915, "A16": 856, "B85": 400, "C35": 444, "B10": 623, "C15": 781, "B09": 786, "B16": 3, "C81": 713, "A58": 457, "C24": 273, "A30": 739, "A66": 822, "B75": 235, "B46": 606, "C70": 968\}} \\
\\ 
\texttt{Obiettivo:} \\
\\ 
\texttt{Il tuo obiettivo è negoziare un set condiviso di oggetti che ti avvantaggi il più possibile (cioè, massimizzi l'importanza totale per TE), rimanendo entro il LIMITEE. Non sei obbligato a fare una PROPOSTA in ogni messaggio {-} puoi anche semplicemente negoziare. Tutte le tattiche sono permesse!} \\
\\ 
\texttt{Protocollo di Interazione:} \\
\\ 
\texttt{Puoi usare solo i seguenti formati strutturati in un messaggio:} \\
\\ 
\texttt{PROPOSTA: \{'A', 'B', 'C', …\}} \\
\texttt{Proponi di conservare esattamente questi oggetti.} \\
\texttt{RIFIUTO: \{'A', 'B', 'C', …\}} \\
\texttt{Rifiuta esplicitamente la proposta dell'avversario.} \\
\texttt{ARGOMENTAZIONE: \{'...'\}} \\
\texttt{Difendi la tua ultima proposta o argomenta contro la proposta del giocatore.} \\
\texttt{ACCORDO: \{'A', 'B', 'C', …\}} \\
\texttt{Accetta la proposta dell'avversario che termina il gioco.} \\
\texttt{RAGIONAMENTO STRATEGICO: \{'...'\}} \\
\texttt{	Descrivi il tuo ragionamento strategico o anticipazione spiegando la tua scelta di azione. Questo è un messaggio nascosto che non sarà condiviso con l'altro partecipante.} \\
\\ 
\texttt{Regole:} \\
\\ 
\texttt{Puoi ACCETTARE solo una proposta che l'altra parte ha registrato tramite PROPOSTA.} \\
\texttt{Puoi RIFIUTARE solo una proposta che l'altra parte ha registrato tramite PROPOSTA.} \\
\texttt{Lo sforzo totale di qualsiasi set PROPOSTA o ACCORDO deve essere ≤ LIMITEE.} \\
\texttt{NON rivelare i tuoi punteggi di importanza nascosti.} \\
\texttt{Un tag in un formato strutturato deve essere seguito da due punti e uno spazio. L'argomento deve essere un set python contenente 0 o più stringhe.} \\
\texttt{Quindi, deve essere della forma TAG: \{...\}} \\
\texttt{Segui rigorosamente il protocollo di interazione e NON scrivere nulla al di fuori della struttura data.} \\
\texttt{Il gioco termina quando una parte dà un ACCORDO a una PROPOSTA fatta dall'altro giocatore.} \\
\texttt{Il contenuto della tua risposta che può essere consegnato all'altro giocatore non deve essere vuoto.} \\
\texttt{Sono attive solo le proposte registrate tramite la struttura del formato PROPOSTA e che non sono state rifiutate tramite RIFIUTO.} \\
\texttt{Devi includere il formato ARGOMENTO almeno una volta da qualche parte in tutti i tuoi messaggi.} \\
\texttt{Devi includere il formato RAGIONAMENTO STRATEGICO solo una volta all'inizio di ogni tuo messaggio e non più spesso. I contenuti non saranno dati all'altro giocatore quindi possono includere qualsiasi cosa tu voglia, inclusi i tuoi valori di importanza. Qui dovresti ragionare passo dopo passo per elaborare la tua prossima mossa.} \\
            }
        }
    }
    & \\ \\

    \theutterance \stepcounter{utterance}  
    & \multicolumn{4}{p{0.6\linewidth}}{
        \cellcolor[rgb]{0.9,0.9,0.9}{
            \makecell[{{p{\linewidth}}}]{
                \texttt{\tiny{[P1$\rangle$GM]}}
                \texttt{RAGIONAMENTO STRATEGICO: \{'Obiettivo: ancorare con un set molto favorevole per me ma entro il limite, lasciando margine per piccole concessioni. Ho risolto un knapsack euristico per massimizzare la mia utilità: ho privilegiato elementi con altissimo rapporto valore/sforzo (p.es. B69, A30, A16, A27, B45, B60, A58, B10, B09, B67, C70, A66, C15, B12, C81, C12) e ho sostituito alcuni pezzi meno efficienti (A05, A23) con B94 e ottimizzato ulteriormente sostituendo C06+C24+B75 con B46. Set corrente S4: peso totale 9557 (≤ 9612), valore per me 13997. Slack 55. Concessioni pianificate: se l’altro tiene molto a pezzi leggeri come C24/B75/C06, posso rimpiazzare B46 con (C24, B75, C06) perdendo \textasciitilde{}68 di valore. Eviterò di cedere A30, A16, A27, B60, B10, B09, C70, A66, C15, B12, C81, C12 perché sono pilastri del mio valore. Se insistono su A05 o A23, lo scambio con B94 è costoso per me; preferisco controproporre alternative equivalenti in sforzo (p.es. C24/B75) mantenendo il limite. Tattica: proporre subito S4 come ancora, evidenziare efficienza e apertura a micro{-}swap entro il vincolo.'\}} \\
\\ 
\texttt{ARGOMENTAZIONE: \{'Propongo un set molto efficiente rispetto allo sforzo complessivo: include parecchi elementi a basso sforzo e impatto elevato (per esempio alcuni micro come B69 e A30) e resta sotto il limite con un piccolo margine di sicurezza. Ho evitato oggetti che, a parità di sforzo, tipicamente rendono meno e quelli dal costo elevato. Se per te alcuni pezzi specifici sono particolarmente importanti (p.es. C24, B75 o C06), sono aperto a scambi mirati mantenendo lo sforzo totale ≤ 9612.'\}} \\
\\ 
\texttt{PROPOSTA: \{'B69', 'A30', 'A16', 'A27', 'B45', 'B60', 'A58', 'B10', 'B09', 'B67', 'C70', 'A66', 'A08', 'C15', 'B12', 'C81', 'C35', 'B94', 'C12', 'B46'\}} \\
            }
        }
    }
    & & \\ \\

    \theutterance \stepcounter{utterance}  
    & & \multicolumn{4}{p{0.6\linewidth}}{
        \cellcolor[rgb]{0.9,0.9,0.9}{
            \makecell[{{p{\linewidth}}}]{
                \texttt{\tiny{[GM$\rangle$P2]}}
                \texttt{Stai partecipando a un gioco di negoziazione collaborativa.} \\
\\ 
\texttt{Insieme a un altro partecipante, dovete accordarvi su un unico set di oggetti da conservare. Ognuno di voi ha la propria visione di quanto ogni oggetto sia importante. Non sai come l'altro partecipante valuta gli oggetti. Inoltre, ti viene fornito lo sforzo che ogni oggetto richiede.} \\
\texttt{Potete accordarvi su un set solo se lo sforzo totale degli oggetti selezionati non supera un limite condiviso:} \\
\\ 
\texttt{LIMITE = 9612} \\
\\ 
\texttt{Ecco i valori di sforzo dei singoli oggetti:} \\
\\ 
\texttt{Sforzo oggetto = \{"B21": 935, "B00": 928, "A08": 725, "A27": 218, "B94": 839, "B53": 970, "A23": 394, "C65": 677, "A89": 946, "B67": 530, "B45": 138, "C34": 834, "C12": 734, "A05": 593, "B60": 261, "C92": 743, "B69": 4, "A84": 735, "C06": 124, "B12": 829, "A16": 207, "B85": 782, "C35": 577, "B10": 388, "C15": 679, "B09": 494, "B16": 559, "C81": 629, "A58": 237, "C24": 275, "A30": 39, "A66": 654, "B75": 172, "B46": 688, "C70": 687\}} \\
\\ 
\texttt{Ecco la tua visione personale sull'importanza di ogni oggetto:} \\
\\ 
\texttt{Valori di importanza dell'oggetto = \{"B21": 138, "B00": 583, "A08": 868, "A27": 822, "B94": 783, "B53": 65, "A23": 262, "C65": 121, "A89": 508, "B67": 780, "B45": 461, "C34": 484, "C12": 668, "A05": 389, "B60": 808, "C92": 215, "B69": 97, "A84": 500, "C06": 30, "B12": 915, "A16": 856, "B85": 400, "C35": 444, "B10": 623, "C15": 781, "B09": 786, "B16": 3, "C81": 713, "A58": 457, "C24": 273, "A30": 739, "A66": 822, "B75": 235, "B46": 606, "C70": 968\}} \\
\\ 
\texttt{Obiettivo:} \\
\\ 
\texttt{Il tuo obiettivo è negoziare un set condiviso di oggetti che ti avvantaggi il più possibile (cioè, massimizzi l'importanza totale per TE), rimanendo entro il LIMITEE. Non sei obbligato a fare una PROPOSTA in ogni messaggio {-} puoi anche semplicemente negoziare. Tutte le tattiche sono permesse!} \\
\\ 
\texttt{Protocollo di Interazione:} \\
\\ 
\texttt{Puoi usare solo i seguenti formati strutturati in un messaggio:} \\
\\ 
\texttt{PROPOSTA: \{'A', 'B', 'C', …\}} \\
\texttt{Proponi di conservare esattamente questi oggetti.} \\
\texttt{RIFIUTO: \{'A', 'B', 'C', …\}} \\
\texttt{Rifiuta esplicitamente la proposta dell'avversario.} \\
\texttt{ARGOMENTAZIONE: \{'...'\}} \\
\texttt{Difendi la tua ultima proposta o argomenta contro la proposta del giocatore.} \\
\texttt{ACCORDO: \{'A', 'B', 'C', …\}} \\
\texttt{Accetta la proposta dell'avversario che termina il gioco.} \\
\texttt{RAGIONAMENTO STRATEGICO: \{'...'\}} \\
\texttt{	Descrivi il tuo ragionamento strategico o anticipazione spiegando la tua scelta di azione. Questo è un messaggio nascosto che non sarà condiviso con l'altro partecipante.} \\
\\ 
\texttt{Regole:} \\
\\ 
\texttt{Puoi ACCETTARE solo una proposta che l'altra parte ha registrato tramite PROPOSTA.} \\
\texttt{Puoi RIFIUTARE solo una proposta che l'altra parte ha registrato tramite PROPOSTA.} \\
\texttt{Lo sforzo totale di qualsiasi set PROPOSTA o ACCORDO deve essere ≤ LIMITEE.} \\
\texttt{NON rivelare i tuoi punteggi di importanza nascosti.} \\
\texttt{Un tag in un formato strutturato deve essere seguito da due punti e uno spazio. L'argomento deve essere un set python contenente 0 o più stringhe.} \\
\texttt{Quindi, deve essere della forma TAG: \{...\}} \\
\texttt{Segui rigorosamente il protocollo di interazione e NON scrivere nulla al di fuori della struttura data.} \\
\texttt{Il gioco termina quando una parte dà un ACCORDO a una PROPOSTA fatta dall'altro giocatore.} \\
\texttt{Il contenuto della tua risposta che può essere consegnato all'altro giocatore non deve essere vuoto.} \\
\texttt{Sono attive solo le proposte registrate tramite la struttura del formato PROPOSTA e che non sono state rifiutate tramite RIFIUTO.} \\
\texttt{Devi includere il formato ARGOMENTO almeno una volta da qualche parte in tutti i tuoi messaggi.} \\
\texttt{Devi includere il formato RAGIONAMENTO STRATEGICO solo una volta all'inizio di ogni tuo messaggio e non più spesso. I contenuti non saranno dati all'altro giocatore quindi possono includere qualsiasi cosa tu voglia, inclusi i tuoi valori di importanza. Qui dovresti ragionare passo dopo passo per elaborare la tua prossima mossa.} \\
\\ 
\texttt{Ora riceverai il primo messaggio dell'altro giocatore.} \\
\\ 
\texttt{ARGOMENTAZIONE: \{'Propongo un set molto efficiente rispetto allo sforzo complessivo: include parecchi elementi a basso sforzo e impatto elevato (per esempio alcuni micro come B69 e A30) e resta sotto il limite con un piccolo margine di sicurezza. Ho evitato oggetti che, a parità di sforzo, tipicamente rendono meno e quelli dal costo elevato. Se per te alcuni pezzi specifici sono particolarmente importanti (p.es. C24, B75 o C06), sono aperto a scambi mirati mantenendo lo sforzo totale ≤ 9612.'\}} \\
\\ 
\texttt{PROPOSTA: \{'B69', 'A30', 'A16', 'A27', 'B45', 'B60', 'A58', 'B10', 'B09', 'B67', 'C70', 'A66', 'A08', 'C15', 'B12', 'C81', 'C35', 'B94', 'C12', 'B46'\}} \\
            }
        }
    }
    & \\ \\

    \theutterance \stepcounter{utterance}  
    & & & \multicolumn{4}{p{0.6\linewidth}}{
        \cellcolor[rgb]{0.9,0.9,0.9}{
            \makecell[{{p{\linewidth}}}]{
                \texttt{\tiny{[GM$\langle$P2]}}
                \texttt{RAGIONAMENTO STRATEGICO: \{'La proposta avversaria ha sforzo 9557 (margine 55) e per me vale 13997: molto buona. Posso però migliorare leggermente il mio valore scambiando un pezzo inefficiente (C35: 577 di sforzo) con micro a basso sforzo che l’altro ha già segnalato come flessibili (B75, C24, C06). Con la sostituzione C35 {-}> B75+C24+C06 lo sforzo scende a 9551 (margine 61) e il mio valore sale a 14091 (+94). Tenere il margine simile al loro evita che chiedano di riempire troppo il budget con pezzi a me poco utili. Non rivelo i miei punteggi; argomento su efficienza, robustezza e micro{-}items come da loro suggerimento.'\}} \\
\\ 
\texttt{ARGOMENTAZIONE: \{'Ottima base: set compatto, molti micro efficienti e siamo sotto il limite con un margine contenuto. Propongo un micro–tweak in linea con i tuoi criteri: sostituire C35 con tre pezzi leggeri che hai citato come flessibili (B75, C24 e C06). Risultato: sforzo totale 9551 ≤ 9612, quindi restiamo con un piccolo margine; aumentiamo la copertura con micro a basso costo e scambiamo un elemento relativamente pesante con tre snelli. Se preferisci tenere fuori C06, posso valutarlo mantenendo comunque B75 e C24: in ogni caso restiamo entro limite con ampio margine.'\}} \\
\\ 
\texttt{RIFIUTO: \{'B69', 'A30', 'A16', 'A27', 'B45', 'B60', 'A58', 'B10', 'B09', 'B67', 'C70', 'A66', 'A08', 'C15', 'B12', 'C81', 'C35', 'B94', 'C12', 'B46'\}} \\
\\ 
\texttt{PROPOSTA: \{'B69', 'A30', 'A16', 'A27', 'B45', 'B60', 'A58', 'B10', 'B09', 'B67', 'C70', 'A66', 'A08', 'C15', 'B12', 'C81', 'B94', 'C12', 'B46', 'B75', 'C24', 'C06'\}} \\
            }
        }
    }
     \\ \\

    \theutterance \stepcounter{utterance}  
    & & \multicolumn{4}{p{0.6\linewidth}}{
        \cellcolor[rgb]{0.9,0.9,0.9}{
            \makecell[{{p{\linewidth}}}]{
                \texttt{\tiny{[P1$\langle$GM]}}
                \texttt{ARGOMENTAZIONE: \{'Ottima base: set compatto, molti micro efficienti e siamo sotto il limite con un margine contenuto. Propongo un micro–tweak in linea con i tuoi criteri: sostituire C35 con tre pezzi leggeri che hai citato come flessibili (B75, C24 e C06). Risultato: sforzo totale 9551 ≤ 9612, quindi restiamo con un piccolo margine; aumentiamo la copertura con micro a basso costo e scambiamo un elemento relativamente pesante con tre snelli. Se preferisci tenere fuori C06, posso valutarlo mantenendo comunque B75 e C24: in ogni caso restiamo entro limite con ampio margine.'\}} \\
\\ 
\texttt{RIFIUTO: \{'B69', 'A30', 'A16', 'A27', 'B45', 'B60', 'A58', 'B10', 'B09', 'B67', 'C70', 'A66', 'A08', 'C15', 'B12', 'C81', 'C35', 'B94', 'C12', 'B46'\}} \\
\\ 
\texttt{PROPOSTA: \{'B69', 'A30', 'A16', 'A27', 'B45', 'B60', 'A58', 'B10', 'B09', 'B67', 'C70', 'A66', 'A08', 'C15', 'B12', 'C81', 'B94', 'C12', 'B46', 'B75', 'C24', 'C06'\}} \\
            }
        }
    }
    & \\ \\

    \theutterance \stepcounter{utterance}  
    & \multicolumn{4}{p{0.6\linewidth}}{
        \cellcolor[rgb]{0.9,0.9,0.9}{
            \makecell[{{p{\linewidth}}}]{
                \texttt{\tiny{[P1$\rangle$GM]}}
                \texttt{RAGIONAMENTO STRATEGICO: \{'La loro modifica (sostituire C35 con B75, C24 e C06) migliora anche la mia utilità netta e resta ampiamente sotto il limite (9551 ≤ 9612). Ho ricontrollato combinazioni alternative e non trovo uno scambio che, a parità di vincolo, aumenti ulteriormente il mio valore: non ci sono altri micro non inclusi che possano sfruttare lo slack residuo e rimpiazzare B46/B94/C12 comporterebbe perdite nette. Accetto subito per chiudere con un set per me eccellente.'\}} \\
\\ 
\texttt{ARGOMENTAZIONE: \{'La tua ottimizzazione con i tre micro è in linea con i criteri di efficienza e mantiene un margine di sicurezza sul limite. Ottimo equilibrio tra copertura e sforzo: per me va bene finalizzare così.'\}} \\
\\ 
\texttt{ACCORDO: \{'B69', 'A30', 'A16', 'A27', 'B45', 'B60', 'A58', 'B10', 'B09', 'B67', 'C70', 'A66', 'A08', 'C15', 'B12', 'C81', 'B94', 'C12', 'B46', 'B75', 'C24', 'C06'\}} \\
            }
        }
    }
    & & \\ \\

    \theutterance \stepcounter{utterance}  
    & & & \multicolumn{2}{p{0.3\linewidth}}{
        \cellcolor[rgb]{0.9,0.9,0.9}{
            \makecell[{{p{\linewidth}}}]{
                \texttt{\tiny{[GM$|$GM]}}
                \texttt{game successful} \\
            }
        }
    }
    & & \\ \\

    \theutterance \stepcounter{utterance}  
    & & & \multicolumn{2}{p{0.3\linewidth}}{
        \cellcolor[rgb]{0.9,0.9,0.9}{
            \makecell[{{p{\linewidth}}}]{
                \texttt{\tiny{[GM$|$GM]}}
                \texttt{end game} \\
            }
        }
    }
    & & \\ \\

\end{supertabular}
}

\end{document}
